Dins de la nucleosíntesi, el procés S es dona a les supernoves i és el responsable d'aproximadament la meitat dels elements més pesants que el ferro. Al final del procés S apareix un cicle tancat de bismut. La primera part d'aquest cicle correspon a la cadena de reaccions següent: el bismut ${}^{209}_{83}\mathrm{Bi}$ captura un neutró i cedeix rajos gamma. El producte d'aquesta primera reacció es desintegra en poloni seguint un decaïment $\beta^-$. En la tercera desintegració, el nucli de poloni es desintegra en un nucli de plom i una partícula $\alpha$.

\begin{apartat}{1,25}
Escriviu les tres reaccions que componen aquesta primera part del cicle.
\end{apartat}

\begin{apartat}{1,25}
La part final, que tanca el cicle, es compon de dues reaccions. En la primera, el nucli de plom captura 3 neutrons, i, en la segona, s'alliberen 1 electró i 1 antineutrí, de manera que s'obté el nucli inicial de bismut ${}^{209}_{83}\mathrm{Bi}$. Escriviu aquestes dues reaccions finals. Escriviu també la reacció resultant del balanç del cicle tancat del bismut, és a dir, el balanç final dels dos cicles, que correspon a la reacció dels 4 neutrons que entren en el cicle.
\end{apartat}
